\documentclass[11pt,a4paper]{article}
\usepackage[T1]{fontenc}
\usepackage{lmodern,amsmath,amssymb,amsthm,mathtools,booktabs,microtype}
\usepackage[margin=26mm]{geometry}
\usepackage[colorlinks=true,linkcolor=blue,urlcolor=blue,citecolor=blue]{hyperref}
\newtheorem{theorem}{Theorem}[section]
\newtheorem{lemma}[theorem]{Lemma}
\newtheorem{proposition}[theorem]{Proposition}
\newtheorem{corollary}[theorem]{Corollary}
\theoremstyle{remark}\newtheorem{remark}[theorem]{Remark}
\newcommand{\Q}{\mathbb Q}\newcommand{\Z}{\mathbb Z}
\newcommand{\spanQ}{\operatorname{span}_{\mathbb Q}}
\newcommand{\fracpart}[1]{\{#1\}}
\title{Sondow Integrals and Three-Grid Cantor Series\\[3pt]
\large Rigidity, Logarithm Elimination, and Remaining Obstructions}
\author{Research working note prepared for Keio}
\date{30 October 2026}
\begin{document}
\maketitle
\begin{abstract}
We record an audited continuation of an attempt to prove the irrationality of
Euler's constant using Sondow's integrals and a geometric Cantor series with
bases $2,4,5$. A rational recurrence for the integer coefficient of $\gamma$
gives a polynomial fractional-part criterion equivalent to irrationality. A
family of geometric series consequently has limit $2$ in the rational case and
lower limit $4/3$ in the irrational case. Further results concern the exact
prime-logarithm span, the absence of a fixed-order rational-function recurrence
for the integrals, prime-square power structure, and pairwise transcendental
differences of normalized integrals. These are auxiliary results: this note
does not establish which case holds for $\gamma$. Priority for these deductions
has not been established by an exhaustive literature review.
\end{abstract}

\section{Definitions and established input}
Throughout, $\log$ is the real natural logarithm, $n\ge1$, and
\[
 d_n=\operatorname{lcm}(1,\ldots,n),\qquad
 B_n=\binom{2n}{n},\qquad H_j=\sum_{r=1}^j\frac1r,\quad H_0=0.
\]
Sondow's integral and its decomposition are \cite{sondow}
\begin{align}
 I_n&=\int_0^1\!\int_0^1
 \frac{[x(1-x)y(1-y)]^n}{(1-xy)(-\log(xy))}\,dx\,dy,\label{integral}\\
 I_n&=B_n\gamma+L_n-A_n,\qquad
 A_n=\sum_{i=0}^n\binom ni^2H_{n+i},\label{sondow-id}\\
 L_n&=2\sum_{k=1}^n\sum_{i=0}^{\min(k-1,n-k)}
 \binom ni^2(H_{n-i}-H_i)\log(n+k).\label{log-form}
\end{align}
In particular, $A_n\in\Q$ and $L_n$ is a positive rational combination of
$\log(n+1),\ldots,\log(2n)$. Define integers $S_n>1$, $C_n$, and $M_n$ by
\[
 \log S_n=d_{2n}L_n,\qquad C_n=d_{2n}B_n,\qquad M_n=d_{2n}A_n\in\Z.
\]
The integer product defining $S_n$ is
\[
 S_n=\prod_{k=1}^n\prod_{i=0}^{\min(k-1,n-k)}
       \prod_{j=i+1}^{n-i}(n+k)^{(2d_{2n}/j)\binom ni^2}.
\]
With $E_n=d_{2n}I_n$ and $b_n=\fracpart{\log S_n}$, the correct sign convention is
\begin{equation}
 E_n=C_n\gamma+\log S_n-M_n.\label{scaled-id}
\end{equation}
We use Sondow's elementary estimates
\begin{equation}
 0<I_n<16^{-n},\qquad d_{2n}<8^n,\qquad 0<E_n<2^{-n}.\label{bounds}
\end{equation}
If $\gamma=u/q\in\Q$, then $q\mid d_{2n}$ for all sufficiently large $n$.
Equation \eqref{scaled-id} and $0<E_n<1$ then give
\begin{equation}
 b_n=E_n\quad\hbox{for all sufficiently large }n.\label{rational-case}
\end{equation}

\section{A polynomial fractional-part criterion}
\begin{lemma}\label{lcm-ratio}
Set $R_n=d_{2n+2}/d_{2n}$. Then $R_n\in\Z_{>0}$,
\[
 R_n\le4n+2,\qquad (n+1)C_{n+1}=(4n+2)R_nC_n.
\]
\end{lemma}
\begin{proof}
The least common multiple increases at an integer $m$ only when $m$ is a prime
power, in which case its ratio is that prime. The new odd integer $2n+1$
contributes at most $2n+1$; the new even integer $2n+2$ contributes at most $2$.
This proves the bound. The identity follows from
$B_{n+1}/B_n=(4n+2)/(n+1)$.
\end{proof}

\begin{theorem}\label{poly-criterion}
Let $\delta_n=1/[16(n+1)^2]$. Then
\begin{equation}
 \gamma\notin\Q
 \quad\Longleftrightarrow\quad
 b_n\ge\delta_n\ \hbox{for infinitely many }n.\label{criterion}
\end{equation}
Consequently, Sondow's sufficient condition $b_n\ge2^{-n}$ infinitely often
is also necessary for irrationality.
\end{theorem}
\begin{proof}
Put
\[
 k_n=M_n-\lfloor\log S_n\rfloor\in\Z,\qquad
 \varepsilon_n=C_n\gamma-k_n=E_n-b_n.
\]
Suppose $b_n<\delta_n$ eventually. Since $2^{-n}<\delta_n$ for $n\ge12$,
we have $|\varepsilon_n|<\delta_n$ eventually. Lemma~\ref{lcm-ratio} gives an integer
\[
 D_n=(n+1)k_{n+1}-(4n+2)R_nk_n
     =(4n+2)R_n\varepsilon_n-(n+1)\varepsilon_{n+1}
\]
with
\begin{align*}
 |D_n|&<\frac{(4n+2)^2}{16(n+1)^2}
                +\frac{n+1}{16(n+2)^2}\\
 &<1-\frac{15}{16(n+1)}+\frac1{4(n+1)^2}<1.
\end{align*}
Thus $D_n=0$ eventually. Comparing with the recurrence for $C_n$ yields
$k_{n+1}/C_{n+1}=k_n/C_n$ eventually. These ratios are one fixed rational
number and converge to $\gamma$, proving $\gamma\in\Q$.

Conversely, \eqref{rational-case} implies $b_n<\delta_n$ eventually if
$\gamma\in\Q$. This proves \eqref{criterion}. Since
$\delta_n>2^{-n}$ eventually, irrationality implies Sondow's exponential
condition infinitely often; the reverse implication follows from
\eqref{rational-case}.
\end{proof}

\begin{corollary}
If $\gamma$ is rational then $n^2b_n\to0$. If $\gamma$ is irrational then
$\limsup_{n\to\infty}n^2b_n\ge1/16$.
\end{corollary}
\begin{remark}
This is a rigorously proved dichotomy, not a proof of the infinite-occurrence
assertion for the actual sequence $S_n$. The recurrence argument becomes
applicable only after assuming eventual smallness.
\end{remark}

\section{Consequences for the three-grid series}
For a positive direction $(a,b,c)$, follow the ray $t(a,b,c)$ from the origin
for $t>0$. At each positive-integer coordinate plane crossing append $2$, $4$,
or $5$ according as the coordinate is $x$, $y$, or $z$. Process simultaneous
crossings in the fixed order $x,y,z$. If the resulting sequence is $(q_j)$, set
\[
 \Omega(a,b,c)=1+\sum_{j\ge1}\frac1{q_1\cdots q_j}.
\]
For any normalized tail of this series, $1/4\le T\le1$. In particular,
after an initial block of $m$ symbols $2$,
$|\Omega-2|\le 2^{-m}$; after an initial block of $m$ symbols $4$,
$|\Omega-4/3|\le(2/3)4^{-m}$.

\begin{theorem}\label{grid-limit}
Set
\[
 a_n=\frac1{16(n+1)^3},\qquad U_n=\Omega(a_n,b_n,E_n).
\]
Then
\[
 \begin{cases}
 U_n\longrightarrow2,&\gamma\in\Q,\\[2pt]
 \displaystyle\liminf_{n\to\infty}U_n=4/3,&\gamma\notin\Q.
 \end{cases}
\]
\end{theorem}
\begin{proof}
By \eqref{bounds}, $E_n/a_n<16(n+1)^3 2^{-n}\to0$. Before the first $z$ crossing,
there are at least $\lfloor a_n/E_n\rfloor$ $x$ crossings. Hence the initial
block avoiding the symbol $5$ has length tending to infinity. All its symbols
are $2$ or $4$, so comparison with an all-$4$ block gives
$\liminf U_n\ge4/3$.

If $\gamma\notin\Q$, Theorem~\ref{poly-criterion} supplies infinitely many $n$
with $b_n\ge\delta_n$, hence $b_n/a_n\ge n+1$. Along these indices,
$E_n<a_n$ eventually and the number of initial $y$ crossings tends to infinity.
The initial block consists entirely of $4$'s, forcing $U_n\to4/3$ on that
subsequence.

If $\gamma\in\Q$, eventually $b_n=E_n$. The number of initial $x$ crossings
is then at least $\lfloor a_n/E_n\rfloor\to\infty$, so $U_n\to2$.
The stated tie convention does not affect any of these limits.
\end{proof}

There is also a fixed-gap formulation. For $n\ge12$ define
$V_n=\Omega(\delta_n,b_n,E_n)$. The first symbol is $4$ when $b_n>\delta_n$,
and $2$ when $b_n<\delta_n$, because $E_n<\delta_n$.
Equality cannot occur: $\log S_n$, hence $b_n$, is transcendental by the
Hermite--Lindemann theorem. Therefore
\[
 b_n\ge\delta_n\iff V_n\le\frac32,
 \qquad b_n<\delta_n\Longrightarrow V_n\ge\frac{13}{8}.
\]
Theorem~\ref{poly-criterion} makes $V_n\le3/2$ infinitely often an equivalent
criterion for irrationality. Neither the separation of these ranges nor the
limit dichotomy determines the actual branch.

\section{What the prime factors prove, and what they do not}
\begin{proposition}\label{prime-power}
For every odd prime $p$, there is an integer $T_p>0$ such that
\[
 S_{p-1}=T_p^{p^2},\qquad v_p(C_{p-1})=2,\qquad
 M_{p-1}\equiv-\frac{d_{2p-2}}p\not\equiv0\pmod p.
\]
\end{proposition}
\begin{proof}
The exponent of the displayed base $n+k$ in the product for $S_n$ is
\[
 e_{n,k}=2d_{2n}\sum_{i=0}^{\min(k-1,n-k)}
             \binom ni^2(H_{n-i}-H_i).
\]
For $n=p-1$, the summands in
$H_{p-1-i}-H_i=\sum_{j=i+1}^{p-1-i}1/j$ pair under $j\mapsto p-j$.
Each pair is $p/[j(p-j)]$, whose denominator is prime to $p$.
Since $v_p(d_{2p-2})=1$, every $e_{p-1,k}$ is divisible by $p^2$.
Also $v_p\binom{2p-2}{p-1}=1$, proving the assertion about $C$.

Finally, $\binom{p-1}{i}^2\equiv1\pmod p$. In $d_{2p-2}H_{p-1+i}$,
only the term of denominator $p$ can contribute modulo $p$, and it occurs
exactly for $i\ge1$. Summing gives
$M_{p-1}\equiv(p-1)d_{2p-2}/p\pmod p$.
\end{proof}

These congruences do not contradict a small fractional part. Indeed, for any
$r\in\{0,\ldots,p^2-1\}$ and $0<\eta<1$, density of
$\{j\log2\}$ supplies $j>0$ with
\[
 \frac r{p^2}<\{j\log2\}<\frac{r+\eta}{p^2}.
\]
The integer $N=2^{p^2j}$ then has
$\lfloor\log N\rfloor\equiv r\pmod{p^2}$ and
$0<\{\log N\}<\eta$. This is a counterexample to an inference from power
structure and congruences alone, not a counterexample involving Sondow's
specific exponents.

\section{The exact logarithm space and the recurrence obstruction}
\begin{theorem}\label{log-rank}
For every $N\ge1$,
\[
 \spanQ\{L_1,\ldots,L_N\}
 =\spanQ\{\log p:p\le2N,\ p\text{ prime}\}.
\]
Its dimension is $\pi(2N)$.
\end{theorem}
\begin{proof}
Inclusion from left to right follows by factoring the integers in
\eqref{log-form}. The prime logarithms are linearly independent over $\Q$
by unique factorization. For the reverse inclusion, $L_1=2\log2$.
For an odd prime $p\le2N$, put $m=(p+1)/2\le N$. In $L_m$, the coefficient
of $\log p$ is positive. All other primes appearing are less than $p$:
the possible bases range from $m+1$ to $2m=p+1$, and the even integer
$p+1$ has no prime factor exceeding $p$. Induction over the primes, or a
triangular change of basis, proves the claim.
\end{proof}

\begin{corollary}\label{not-recursive}
The sequence $(I_n)$ satisfies no nonzero fixed-order homogeneous linear
recurrence with coefficients in $\Q(n)$ valid for all sufficiently large $n$.
\end{corollary}
\begin{proof}
By \eqref{sondow-id} and Theorem~\ref{log-rank},
\[
 \dim_{\Q}\spanQ\{I_1,\ldots,I_N\}\ge\pi(2N)-2\longrightarrow\infty.
\]
A fixed-order recurrence over $\Q(n)$, with leading coefficient not identically
zero, expresses all sufficiently late terms as rational linear combinations
of finitely many earlier terms. This would bound that dimension independently
of $N$, a contradiction.
\end{proof}

This obstruction concerns the particular sequence \eqref{integral} and exact
recurrences of fixed order. It does not rule out growing-order elimination,
recurrences for different approximants, or nonrational coefficients.

\section{An unconditional arithmetic consequence for the integrals}
\begin{theorem}\label{pair-trans}
For distinct positive integers $m,n$, the number
\[
 \frac{I_m}{B_m}-\frac{I_n}{B_n}
\]
is transcendental. Consequently, among all $I_n$ there is at most one
algebraic value. The same at-most-one assertion holds for $E_n=d_{2n}I_n$.
\end{theorem}
\begin{proof}
Write $\alpha_n=A_n/B_n$, $F_n=L_n/B_n$, and $J_n=I_n/B_n$.
Let $K_n$ have probabilities $w_n(i)=\binom ni^2/B_n$, for $0\le i\le n$.
Vandermonde's identity makes these a probability distribution, and
$\alpha_n=\mathbb E H_{n+K_n}$.
For $0\le i\le n$,
\[
 \frac{w_{n+1}(i)}{w_n(i)}
  =\frac{B_n}{B_{n+1}}\frac{(n+1)^2}{(n+1-i)^2}
\]
increases with $i$, and the new distribution also has positive mass at $n+1$.
The ratios cross $1$ at most once, so the cumulative probabilities for
$K_{n+1}$ do not exceed those for $K_n$. Thus $K_{n+1}$ stochastically
dominates $K_n$, and
\[
 \alpha_{n+1}-\alpha_n
 \ge\mathbb E\frac1{n+1+K_n}\ge\frac1{2n+1}.
\]
Since $F_n=\alpha_n-\gamma+J_n$, estimate \eqref{bounds} gives
\[
 F_{n+1}-F_n>
 \frac1{2n+1}-\frac{16^{-n}}{B_n}>0.
\]
For $m>n$, $F_m-F_n$ is therefore a nonzero rational linear combination
of logarithms of positive integers. Its exponential is a positive algebraic
number different from $1$. The Hermite--Lindemann theorem implies that
$F_m-F_n$ is transcendental \cite{baker}. Subtracting the rational number
$\alpha_m-\alpha_n$ proves the assertion about $J_m-J_n$.
Two algebraic values of $I$ would give an algebraic difference, a contradiction.
Multiplication by the nonzero integer $d_{2n}$ preserves algebraicity.
\end{proof}

\begin{remark}
Theorem~\ref{pair-trans} does not distinguish rational from irrational
$\gamma$. If $\gamma$ were algebraic, every $I_n$ would already be
transcendental by \eqref{sondow-id}, since $L_n$ is a nonzero logarithm of an
algebraic number. The theorem is an auxiliary conclusion about the integrals.
\end{remark}

\section{Exact elimination experiments and the unresolved step}
Exact integer arithmetic gives the concrete cancellation
\begin{align*}
 &-3923225775E_1-158150125E_2+392805E_3\\
 &\hspace{17mm}-200868E_4+35150E_5
 =-16097728500\gamma+8969752951.
\end{align*}
All prime-logarithm coefficients cancel exactly. The displayed expression
has magnitude approximately $3.22\times10^8$; this combination does not
produce a small integer linear form. Integer-kernel/LLL probes in the
windows $[10,30]$, $[30,60]$, and $[50,100]$ likewise produced large forms.
These computations test specific reductions, not all possible combinations
and not an asymptotic impossibility theorem.

For a growing block of indices, let $\mathcal L$ be the integer matrix of
prime exponents in $\log S_n$. A vector $u$ in its integer kernel gives
\[
 \sum_nu_nE_n=\left(\sum_nu_nC_n\right)\gamma-\sum_nu_nM_n.
\]
The required assertions are separate:
\begin{enumerate}
\item the coefficients must be small enough that $\sum_n|u_n|E_n\to0$;
\item the resulting form must not vanish identically (or two coefficient
      pairs must be demonstrably independent).
\end{enumerate}
A nonzero kernel vector need not satisfy the second condition: it can also
annihilate the $C$ and $M$ rows. Siegel's lemma alone does not prevent this.
The positivity of \eqref{integral} proves $\int P^2\,d\mu>0$ for a nonzero
polynomial $P$, but squaring a log-annihilating polynomial need not preserve
logarithm cancellation. No estimate resolving these two issues simultaneously
was obtained in this attempt.

\paragraph{Current status.}
The polynomial criterion, geometric dichotomy, prime-power statements,
logarithm-span theorem, recurrence obstruction, and normalized-integral
transcendence theorem above have proofs. The unconditional assertion
$b_n\ge1/[16(n+1)^2]$ infinitely often remains unproved here.
Accordingly, neither the irrationality nor the transcendence of $\gamma$
has been established by this work.

\appendix
\section{The exact tail gcd of the integer coefficients}
\begin{proposition}
For $N\ge1$,
\[
 \gcd_{n\ge N}C_n=
 \begin{cases}
 2d_{2N},&N\text{ is a power of }2,\\
 4d_{2N},&\text{otherwise}.
 \end{cases}
\]
The same gcd is obtained from the finite block $N\le n\le2N^2$.
\end{proposition}
\begin{proof}
Fix an odd prime $p\le2N$ and $k=\lfloor\log_p(2N)\rfloor$.
The lower bound $v_p(C_n)\ge k$ holds for $n\ge N$. Take
$m=(p^{k+1}-1)/2\ge N$. Its base-$p$ digits are all $(p-1)/2$,
so doubling creates no carries and $v_p(B_m)=0$ by the elementary
factorial valuation formula. Also $v_p(d_{2m})=k$.
Thus the gcd has exactly the odd-prime valuations of $d_{2N}$.
Primes greater than $2N$ cannot divide the gcd because they do not divide $C_N$.

For $p=2$, the same factorial valuation formula gives
\[
 v_2(C_m)=\lfloor\log_2(2m)\rfloor+s_2(m),
\]
where $s_2(m)$ is the number of binary ones. Set $k=\lfloor\log_2(2N)\rfloor$.
If $N=2^{k-1}$ the minimum is $k+1$, attained at $m=N$.
Otherwise every $N\le m<2^k$ has at least two binary ones; for $m\ge2^k$
the minimum is also at least $k+2$, attained at $m=2^k$.
This proves the formula. The odd-prime witnesses satisfy
$m<pN\le2N^2$, and the binary witness is at most $2N$.
\end{proof}
This narrows the possible denominator if an eventual small-tail assumption
has a specified starting index. It does not exclude all denominators when
that starting index is unknown.

\begin{thebibliography}{9}
\bibitem{sondow}
J. Sondow, \emph{Criteria for irrationality of Euler's constant},
Proc. Amer. Math. Soc. \textbf{131} (2003), 3335--3344.
\href{https://doi.org/10.1090/S0002-9939-03-07081-3}{doi:10.1090/S0002-9939-03-07081-3}.
Preprint: \url{https://arxiv.org/abs/math/0209070}.
\bibitem{baker}
A. Baker, \emph{Transcendental Number Theory}, Cambridge University Press, 1975.
Classical background on exponential and logarithmic transcendence.
\end{thebibliography}
\end{document}
